% Irodov Part II section 2.4, original 113--121, printed p.76, PDF p.85.
% Carnot, Otto, Brayton and mixed reversible ideal-gas engine cycles.
\thermoentropyproblem{113}{76}
在温度分别为 $T_h>T_c>0$ 的两个热源间工作的
可逆卡诺热机，哪种方式对提高热效率更有效：
将高温热源的温度升高相同的
$\Delta T>0$，
还是将低温热源的温度降低 $\Delta T$
（假设 $T_c>\Delta T$）？
\exsolution
卡诺效率
$\eta=1-T_c/T_h$。
提高高温热源温度得到
\[
\eta_h=1-\frac{T_c}{T_h+\Delta T},
\]
降低低温热源温度得到
\[
\eta_c=1-\frac{T_c-\Delta T}{T_h}.
\]
直接相减：
\[
\eta_c-\eta_h
=\frac{\Delta T}{T_h}
-\frac{T_c\Delta T}{T_h(T_h+\Delta T)}
=\frac{\Delta T(T_h+\Delta T-T_c)}
{T_h(T_h+\Delta T)}>0.
\]
故
\[
\boxed{\text{降低低温热源温度的同量 }\Delta T
\text{ 可得到更大的效率提高}}.
\]
\excheck 比较必须在仍有
$T_c-\Delta T>0$ 的物理范围中进行。

\thermoentropyproblem{114}{76}
氢气作为卡诺热机的工质。
若绝热膨胀阶段：
(a) 气体体积增加 $n=2.0$ 倍；
(b) 气体压强下降 $n=2.0$ 倍，
求各情况下的卡诺效率。
\exsolution
卡诺循环两等温温度为 $T_h,T_c$，
效率
$\eta=1-T_c/T_h$。
假设氢气为刚性双原子理想气体，
$\gamma=7/5$。
(a) 可逆绝热过程满足
$TV^{\gamma-1}=\text{常数}$，
若 $V_c/V_h=n$，则
$T_c/T_h=n^{-(\gamma-1)}$。
因此
\[
\boxed{\eta_a=1-n^{1-\gamma}
=1-2^{-2/5}\simeq0.24\approx25\%}.
\]
(b) 由可逆绝热关系
$T\propto p^{(\gamma-1)/\gamma}$，
$T_c/T_h=n^{-(\gamma-1)/\gamma}$。
故
\[
\boxed{\eta_b=1-n^{-(\gamma-1)/\gamma}
=1-2^{-2/7}\simeq0.18}.
\]
\excheck 两问给定的是不同的绝热压缩比，
不能把体积倍率与压力倍率直接混用。

\thermoentropyproblem{115}{76}
一台卡诺热机原先工作在两个热源之间，
其热机效率为 $\eta=10\%$。
若将同一装置倒转用作卡诺制冷机，
且冷热热源仍不变，
求制冷系数 $\varepsilon$。
\exsolution
对可逆卡诺循环，
热机效率
$\eta=1-T_c/T_h$。
反向制冷机吸收冷源热量 $Q_c$、
消耗外功 $W=Q_h-Q_c$，
制冷系数定义为
$\varepsilon=Q_c/W$。
卡诺关系
$Q_c/Q_h=T_c/T_h$，
因此
\[
\varepsilon=\frac{T_c}{T_h-T_c}
=\frac{T_c/T_h}{1-T_c/T_h}
=\frac{1-\eta}{\eta}.
\]
代入 $\eta=0.10$：
\[
\boxed{\varepsilon=\frac{0.90}{0.10}=9.0}.
\]
\excheck 制冷系数不是热机效率，
可以大于1，因为它衡量的是
搬运的冷源热量与外功之比。

\thermoentropyproblem{116}{76}
一理想气体经过由等温线和绝热线交替组成的
可逆热机循环（原图2.2）。
三段等温过程分别处于
$T_1>T_2>T_3$。
其中两个等温膨胀阶段的体积均以相同倍率
$r>1$ 增大。
求整个循环效率。
\exsolution
在温度 $T_1,T_2$ 上，
各发生一次等温膨胀，
理想气体等温吸热分别为
$Q_1=\nu RT_1\ln r$、
$Q_2=\nu RT_2\ln r$。
绝热连接段不交换热。
由于闭合循环内总熵变化为零，
低温 $T_3$ 的等温压缩须排出
$Q_3$ 满足
\[
\frac{Q_1}{T_1}+\frac{Q_2}{T_2}
-\frac{Q_3}{T_3}=0.
\]
故
\[
Q_3=2\nu RT_3\ln r.
\]
热机做净功
$W=Q_1+Q_2-Q_3$，
热效率为
\[
\boxed{\eta=\frac{W}{Q_1+Q_2}
=1-\frac{2T_3}{T_1+T_2}}.
\]
\excheck 倍率 $r$ 通过等温吸热表达式相消；
若忽略第二条中温等温线，
会错误地套用普通两热源卡诺公式。

\thermoentropyproblem{117}{76}
一理想热机循环由两条等容线和
两条可逆绝热线构成（Otto 循环）。
若体积比 $r=V_{\max}/V_{\min}=10$，
且工质为氮气，求效率。
\exsolution
在绝热压缩与绝热膨胀上，
$TV^{\gamma-1}=\text{常数}$。
令四状态的两等容吸放热
$Q_{\rm in}=\nu C_V(T_3-T_2)$、
$Q_{\rm out}=\nu C_V(T_4-T_1)$。
配对使用绝热关系
$T_2/T_1=T_3/T_4=r^{\gamma-1}$，
可得
\[
\frac{Q_{\rm out}}{Q_{\rm in}}
=r^{1-\gamma}.
\]
故循环效率
\[
\boxed{\eta=1-r^{1-\gamma}}.
\]
氮气取 $\gamma=7/5$，代入 $r=10$：
\[
\boxed{\eta=1-10^{-2/5}\simeq0.60=60\%}.
\]
\excheck 与两条\emph{等压}的循环不同，
这里的压缩倍率定义为体积倍率。

\thermoentropyproblem{118}{76}
一理想气体的可逆热机循环
由两条等压线和两条绝热线组成
（Brayton 循环）。
若循环中最高与最低压强之比为 $n$，
气体绝热指数为 $\gamma$，
求热效率。
\exsolution
设高低压强分别为
$p_h,p_l$，$p_h/p_l=n>1$。
在两绝热支路上
$T\propto p^{(\gamma-1)/\gamma}$，
故成对状态的温度比
\[
\frac{T_2}{T_1}=\frac{T_3}{T_4}
=n^{(\gamma-1)/\gamma}.
\]
高压等压支路吸热
$Q_{\rm in}=\nu C_p(T_3-T_2)$，
低压等压支路排热
$Q_{\rm out}=\nu C_p(T_4-T_1)$。
因此
\[
\frac{Q_{\rm out}}{Q_{\rm in}}
=n^{-(\gamma-1)/\gamma}.
\]
效率为
\[
\boxed{\eta=1-n^{-(\gamma-1)/\gamma}}.
\]
\excheck 必须保证循环方向为正功热机，
并使用正的吸热、排热量计算效率。

\thermoentropyproblem{119}{76}
一理想气体的顺时针循环由两条等容线
与两条等压线组成，
气体绝热指数为 $\gamma$。
若等容加热和等压膨胀时
气体的绝对温度都各自增加 $n$ 倍，
求循环效率。
\exsolution
取初态
$(p_1,V_1,T_1)$，
等容加热使压强变为 $np_1$、
温度变为 $nT_1$；
随后在高压 $np_1$ 下等压膨胀，
温度再增至 $n^2T_1$、
体积变为 $nV_1$。
剩下两支路为等容冷却和低压等压压缩。
循环矩形面积是净功：
\[
W=(np_1-p_1)(nV_1-V_1)
=p_1V_1(n-1)^2
=\nu RT_1(n-1)^2.
\]
吸热在等容加热和高压等压膨胀段：
\[
Q_{\rm in}=\nu C_VT_1(n-1)
+\nu C_pT_1n(n-1).
\]
因此
\[
\eta=\frac{R(n-1)}{C_V+nC_p}.
\]
用 $C_p=\gamma C_V$、
$R=(\gamma-1)C_V$ 化简：
\[
\boxed{\eta=\frac{(\gamma-1)(n-1)}
{1+\gamma n}}.
\]
\excheck 若 $n\to1$ 循环面积与功趋零，
效率也随之趋零。

\thermoentropyproblem{120}{76}
一理想气体经历由
(a) 等容、绝热、等温三支路；
(b) 等压、绝热、等温三支路组成的可逆热机循环。
假设等温支路始终处于\emph{最低温度}
$T_c$，循环的最高温度与最低温度之比为
$n=T_h/T_c>1$。
求两循环的效率。
\exsolution
(a) 等容加热将温度由 $T_c$ 提升至 $T_h$，
随后可逆绝热膨胀回到 $T_c$，
最后在低温等温压缩回初态。
等容吸热
$Q_{\rm in}=\nu C_V(T_h-T_c)
=\nu C_VT_c(n-1)$。
绝热支路给
$V_{\rm after}/V_{\rm before}
=n^{1/(\gamma-1)}$，
因此低温等温排热
$Q_{\rm out}=\nu RT_c\ln n/(\gamma-1)
=\nu C_VT_c\ln n$。
故
$\eta_a=1-\ln n/(n-1)$。
(b) 等压加热的吸热
$Q_{\rm in}=\nu C_pT_c(n-1)$。
绝热连接使低温等温压缩的体积比为
$n^{\gamma/(\gamma-1)}$，
从而
$Q_{\rm out}=\nu RT_c\gamma\ln n/(\gamma-1)
=\nu C_pT_c\ln n$。
于是两情况均有
\[
\boxed{\eta_a=\eta_b=1-\frac{\ln n}{n-1}}.
\]
\excheck 正循环必须沿吸热支路升温至
$T_h$；若把等温支路误当作最高温，
会得到第121题的不同结果。

\thermoentropyproblem{121}{76}
条件与上一题相同，
但等温支路改在整个循环的
\emph{最高温度} $T_h$ 上进行。
已知最高、最低温比为
$n=T_h/T_c>1$，
求两种循环各自的效率。
\exsolution
(a) 高温等温膨胀吸热，
随后等容冷却至 $T_c$，
再通过绝热压缩回高温起点。
绝热支路给
$V_{\rm high}/V_{\rm low}
=n^{1/(\gamma-1)}$，
高温等温吸热
$Q_{\rm in}=\nu RT_h\ln n/(\gamma-1)
=\nu C_VT_h\ln n$。
等容冷却排热
$Q_{\rm out}=\nu C_V(T_h-T_c)
=\nu C_VT_h(1-1/n)$，
故
$\eta_a=1-(n-1)/(n\ln n)$。
(b) 高温等温膨胀的体积比是
$n^{\gamma/(\gamma-1)}$，
所以
$Q_{\rm in}=\nu C_pT_h\ln n$；
等压冷却排热
$Q_{\rm out}=\nu C_pT_h(1-1/n)$。
两种情况同为
\[
\boxed{\eta_a=\eta_b=1-\frac{n-1}{n\ln n}}.
\]
\excheck $n>1$ 时
$\ln n>(n-1)/n$，
故效率为正且小于1。
